\section{Introduction}
\label{ch:introduction}

Working memory requires a neural system to preserve task-relevant information
after its eliciting input has disappeared. Recurrent excitation, attractor
dynamics, and slow manifolds provide candidate mechanisms for that
maintenance~\cite{wang2013nmda,ghazizadeh2021}. A weak change in those dynamics
need not erase memory at once; it may instead expose selective vulnerability
when information must be held longer, protected from irrelevant input, or
updated under greater load.

Acute human psilocybin studies motivate exactly such selective targets rather
than global cognitive collapse: clearer impairment under higher updating load,
greater vulnerability when competing information must be ignored, stronger
effects at longer intervals or higher dose, and response slowing that can
outpace accuracy loss~\cite{barrett2018,carter2005,wittmann2007,vollenweider1998,yousefi2025}.
Psychedelic theory and modelling further motivate competing computational
interventions---especially gain-like and weighting changes---without specifying
a unique working-memory RNN equation~\cite{carhartharris2019,herzog2023}.
Section~\ref{ch:literature} reviews that evidence in full.

This dissertation therefore uses task-trained recurrent networks as a
controlled testbed. The aim is computational sufficiency and comparative
mechanism: which literature-motivated perturbations, if any, reproduce the
selected behavioural pattern more convincingly than generic disruption, and
how that pattern emerges from recurrent population dynamics
(Section~\ref{ch:aims}). Methods, the completed candidate screen, and the
outstanding specificity and dynamics experiments are developed in
Sections~\ref{ch:methods}--\ref{ch:results-dynamics}.
